\documentclass[11pt,a4paper]{article}

\usepackage[utf8]{inputenc}
\usepackage[T1]{fontenc}
\usepackage[romanian]{babel}
\usepackage{geometry}
\usepackage{booktabs}
\usepackage{hyperref}

\geometry{margin=2.3cm}
\hypersetup{hidelinks}

\title{Raport scurt privind stadiul proiectului de diplomă}
\author{Alin Seicarin}
\date{7 august 2026}

\begin{document}

\maketitle

\section*{Obiectiv și delimitarea temei}

Proiectul urmărește compararea strategiilor PID, LQR și MPC pentru urmărirea
geometrică a unei căi de către un robot mobil cu tracțiune diferențială, simulat
în ROS~2 și Gazebo. Calea nu este parametrizată în timp, iar transferul pe un
robot fizic nu face parte din domeniul lucrării.

\section*{Stadiul actual}

Etapa PID este implementată și stabilizată. Au fost realizate:

\begin{itemize}
    \item modelul robotului și simularea Gazebo, cu odometrie, IMU și fuziune
    senzorială prin EKF;
    \item un manager comun al referinței, bazat pe proiecția continuă a
    robotului pe cale, care furnizează eroarea laterală, eroarea de orientare și
    curbura locală;
    \item un regulator PID în cascadă: bucla exterioară transformă eroarea
    laterală într-o corecție de orientare, iar bucla interioară furnizează
    corecția vitezei unghiulare;
    \item o politică de comandă comună PID--LQR--MPC, care impune aceeași viteză
    de referință, aceeași anticipare a curburii și aceleași saturații;
    \item trasee de test: dreaptă, sinusoidă, viraj, cerc și figură opt;
    \item un test repetabil de respingere a perturbațiilor, prin adăugarea
    temporară a unei erori asupra comenzii de viteză unghiulară;
    \item evaluare independentă față de poziția reală din Gazebo, astfel încât
    erorile de localizare nu sunt confundate cu performanța fizică a
    regulatorului;
    \item scripturi de benchmark și teste automate pentru configurare,
    referință, regulator, perturbare și evaluarea rezultatelor.
\end{itemize}

În verificările preliminare, toate cele cinci trasee au fost finalizate în
toleranța fizică de $0{,}15\,\mathrm{m}$. Valorile RMS ale erorii laterale reale
au fost cuprinse între aproximativ $0{,}002\,\mathrm{m}$ și
$0{,}086\,\mathrm{m}$. La aplicarea unei perturbații de
$0{,}6\,\mathrm{rad/s}$ timp de o secundă, eroarea laterală maximă a fost de
aproximativ $0{,}057\,\mathrm{m}$, revenirea strictă a durat
$4{,}361\,\mathrm{s}$, iar abaterea laterală finală a fost
$0{,}003\,\mathrm{m}$. Aceste rezultate reprezintă validarea funcțională a
implementării, nu încă experimentul comparativ final.

Identificarea de sistem a fost eliminată din planul principal. În simularea
actuală, dinamica acționării este cunoscută și impusă explicit prin modelul
Gazebo și limitarea accelerației roților; pentru LQR și MPC este mai coerentă
folosirea unui model cinematic liniarizat, comun și reproductibil.

\section*{Etape rămase}

\begin{enumerate}
    \item definitivarea exportului și prelucrării datelor în MATLAB: traiectorii,
    eroare laterală, efort de comandă și indicatori de robustețe;
    \item derivarea și discretizarea modelului liniarizat al erorilor pentru
    robotul diferențial;
    \item implementarea și reglarea regulatorului LQR pe aceeași interfață și cu
    aceleași limitări ca PID;
    \item implementarea MPC, inclusiv orizontul de predicție și constrângerile
    asupra comenzilor;
    \item campania experimentală finală: rulări repetate pentru fiecare
    regulator, traseu și perturbare, urmate de analiză statistică;
    \item redactarea rezultatelor, comparației și concluziilor lucrării.
\end{enumerate}

\section*{Concluzie}

Infrastructura comună de simulare și evaluare este funcțională, iar PID-ul în
cascadă oferă o bază de referință validată. Următorul obiectiv tehnic este
implementarea LQR pe modelul analitic liniarizat, urmată de MPC și de comparația
experimentală finală în condiții identice.

\end{document}
